\input{glyphtounicode}
\pdfgentounicode=1
% =============================================================================
% Global Normalization in Joshi's Arithmetic Teichmuller Theory
% Assessment / classification paper -- corrective draft v1.2 (English)
% Project: DRAFT__Global_Normalization_Contract_Lemma  (.RESEARCH/.LAB/.HCT)
% Basis: OUTLINE.md; _proof-notes/LEMMA_FORMULIERUNG_V1_2026-08-12.md;
%        P1_QUELLENPRUEFUNG_NORMIERUNG_2026-08-12.md;
%        P2_TERM_KLASSIFIKATION_THM6101_2026-08-13.md;
%        R-E2_MASSPINNUNG_2026-08-13.md
% v1.1 (13 Aug 2026): branch-conditional revision. Every verdict that depends on
%        the selection rule for lambda_y is now marked for the branch under which
%        it holds (Section 3.5 trilemma); Lemma~1 (rigidity of the product-formula
%        multiplier) added in Section 3.6; Section 7 rewritten (the former
%        three-reading frame is superseded); AI disclosure switched to the
%        pipeline standard by \input. Report: _results/V11_TRILEMMA_REVISION_2026-08-13.md
% This is NOT a proof paper. See the scope box after the abstract.
% =============================================================================

\documentclass[11pt,a4paper]{article}

\usepackage[utf8]{inputenc}
\usepackage{cmap}
\usepackage[T1]{fontenc}
\usepackage{lmodern}
\usepackage{microtype}
\usepackage{xurl}
\usepackage[a4paper,margin=2.7cm]{geometry}
\usepackage{amsmath}
\usepackage{amsthm}
\usepackage{amssymb}
\usepackage{array}
\usepackage{booktabs}
\usepackage[colorlinks=true,linkcolor=black,citecolor=black,urlcolor=blue,unicode=true]{hyperref}
\def\HyperDestNameFilter{en.}


\hyphenation{Ara-ke-lov arith-me-ti-coid arith-me-ti-coids ho-lo-mor-phoid
  ho-lo-mor-phoids nor-mal-i-za-tion pa-ram-e-tri-za-tion char-ac-ter-i-za-tion}

% v1.1: the revision added several paragraphs dense in unbreakable math atoms;
% a little extra stretch keeps them inside the text block.
\emergencystretch=1.5em

% ragged-right paragraph columns for the two classification tables
\newcolumntype{P}[1]{>{\raggedright\arraybackslash}p{#1}}

\hypersetup{
  unicode=true,
  pdfencoding=unicode,
  psdextra,
  breaklinks=true,
  pdftitle={Global Normalization in Joshi's Arithmetic Teichmüller Theory},
  pdfauthor={Lukas Geiger},
  pdfsubject={Assessment of the normalization identification in ATS~IV Theorem~6.10.1},
  pdfkeywords={Arithmetic Teichmüller Theory, Global Normalization, Scale Identification, Invariance Taxonomy, Trilemma, Joshi, Mochizuki, Scholze--Stix}
}

% --- theorem-like environments ------------------------------------------------
\theoremstyle{definition}
\newtheorem{question}{Question}
\theoremstyle{plain}
\newtheorem*{derivationD}{Conditional derivation D-1}
\newtheorem{lemmaR}{Lemma}
\theoremstyle{definition}
\newtheorem*{findingF}{Finding F-NEU-1}
\newtheorem*{criterionK}{Criterion~K (candidate)}

% --- notation -----------------------------------------------------------------
\newcommand{\arith}{\operatorname{arith}}
\newcommand{\hol}{\operatorname{hol}}
\newcommand{\ordd}{\operatorname{ord}}
\newcommand{\Vol}{\operatorname{Vol}}
\newcommand{\LogVol}{\operatorname{LogVol}}
\newcommand{\Thet}{\widetilde{\Theta}}
\newcommand{\Ynum}{Y_L}

% Every source citation is anchored to a named extract, because the four
% extracts carry independent line numbering.
\newcommand{\anchor}[2]{(#1, extract l.~#2)}

\title{Global Normalization in Joshi's Arithmetic Teichmüller Theory:\\
What the Identification Carries --- and What Carries the Theorem}
\author{Lukas Geiger\\[2pt]
\small Independent Researcher, Bernau, Germany\\
\small ORCID: 0009-0005-7296-1534}
\date{2 October 2026 --- Corrective draft v1.2}

\begin{document}
\maketitle

\begin{abstract}
\noindent
This assessment asks what the same global normalization carries in
Theorem~6.10.1 of Joshi's ATS~IV (arXiv:2403.10430v2). Under the explicitly
place-wise Artin--Whaples reading, number-field absolute values are identified
definitionally; the rule selecting the coordinate $\lambda_y$ remains an open
correspondence question. An elementary product-formula lemma constrains
multipliers of the standard valuations to one global scalar. The exclusion of a
fixed product-formula hyperplane remains conditional on assumptions (F1)--(F2).
The three normalization branches are diagnostic, not exhaustive.
This corrective version distinguishes the source's asserted Frobenius stability
of a collated family from the unverified transport to the two indexed hulls in E2.
It corrects three source readings: ATS~IV states its height domain as
$\overline L^*$, E2 equates absolute values of log-volumes, and the algebraic
holomorphic hull of ATS~III is not generically a convex closure. An explicit
counterexample gives $\operatorname{conv}(\mathbb Z_p\cdot1)=\mathbb Z_p\cdot1
\subsetneq\mathcal O_E=\operatorname{hull}(\mathbb Z_p\cdot1)$ in an unramified
quadratic extension $E/\mathbb Q_p$. This does not refute the specialized
tensor-region assertions or the ATS bounds. Normalized additive Haar measure is
kept separate from powers of an absolute value: $m(p\mathbb Z_p)=p^{-1}$ does not
become $p^{-c}$ by changing a valuation exponent. Accordingly the previous
positive claim of cross-frame measure pinning is withdrawn in every branch.
A hull- and measure-compatible indexed transport is named as a possible
sufficient route stronger than the displayed absolute-value equality, not as a
proved or necessary condition. The selection rule, indexed E2 transport,
compatible measure and weight bindings, and order comparison remain open.
No proof or disproof of $abc$, Mochizuki's Corollary~3.12, or Joshi's estimates
is claimed.
\end{abstract}

\clearpage
% -----------------------------------------------------------------------------
% SCOPE AND CLAIM BOUNDARY
% -----------------------------------------------------------------------------
\begin{center}
\fbox{%
\parbox{\dimexpr\textwidth-2\fboxsep-2\fboxrule\relax}{%
\medskip
\textbf{Scope and claim boundary.}\par\medskip
\begin{itemize}\setlength{\itemsep}{2pt}
\item This paper proves \emph{nothing} about the $abc$ conjecture, about
  inter-universal Teichmüller theory, about Mochizuki's Corollary~3.12, or about
  the correctness of Joshi's bounds. It contains no proof and no disproof of any
  of these.
\item It \emph{verifies no bound}. No proof line of ATS~IV Theorem~6.10.1, of
  Corollary~9.11.1.1, or of Corollary~3.12 was followed step by step; no estimate
  was recomputed.
\item What it does: it \emph{classifies definitions and the structure of
  justification} in one named theorem, and states for each of its components
  whether that component is definitional, supplied by construction,
  branch-conditional, or open.
\item It takes no position in the dispute between Joshi and
  Scholze--Stix~\cite{ScholzeStix2018} on the mobility of the product-formula hyperplane. Where
  that dispute is touched, it is reported, not adjudicated.
\item The structural analogy sketched in Section~\ref{sec:b2outlook} is a working
  hypothesis of the present program about the \emph{shape of arguments}. It is
  not a mathematical connection and, in particular, not a claim about the Riemann
  hypothesis.
\item The central open item of this version --- that the passages we read state no
  rule selecting the normalization coordinate of a moved arithmeticoid
  (Section~\ref{sec:trilemma}) --- is a \emph{definitional gap}: a rule we did not
  find, which may well stand in material we have not read. It is not a refutation
  of ATS~III Corollary~4.2.2.4 and not an error claim of any kind. Lemma~1
  (Section~\ref{sec:rigidity}) constrains which rules the product formula
  \emph{alone} can supply; it does not assert that no such rule exists.
\item Where a tension between two passages of a source is recorded, this is an
  invitation to resolve it and a candidate for correspondence with the author ---
  not a claim that it is unresolvable, and not an error claim.
\item \textbf{Textual basis} are \texttt{pdftotext} extracts of the arXiv PDFs.
  All line references (``extract l.~NNNN'') are finding aids into those extracts
  and are \emph{not} pagination of the published documents. Direct visual
  readback covered the source-PDF pages for equation (5.3.3),
  Theorem~5.10.1(3)--(4), Theorem~10.20.1(3),
  Definition~4.4.2/formula (4.4.5), and Theorem~6.10.1. No blanket
  OCR-robustness claim is made; in particular, the E1 absolute-value bars and the
  direction of the fiber exponent were checked in the source PDFs.
\end{itemize}
\medskip
}}
\end{center}

\clearpage
\tableofcontents
\clearpage

% =============================================================================
\section{Introduction: the point at issue}
% =============================================================================

\subsection{The finding that prompted this paper}

The present paper is the fourth part of a small research program whose third
part was a structured ``ledger test'' of the argumentative architecture of ATS~IV:
a requirement-by-requirement check of which load-bearing steps of the paper are
supplied, which are referenced, and which are asserted. That test returned one
requirement --- the one concerning the comparability of the two arithmeticoids
entering Theorem~6.10.1 --- as met only \emph{in part}. The reason was not a gap
in a computation. It was that the identification carrying the middle equality of
the theorem is, in the text's own word, a tautology
\anchor{ATS~IV}{3272--3273}, while the very same theorem's remark forbids the
place-by-place comparison that a reader might otherwise use to check it.

The two textual poles are these. On the one hand the theorem opens with a choice:

\begin{quote}\small
``Choose the same global normalization for the number field provided by the
arithmeticoids $y_0'$ and $y_0$. With this choice of simultaneous global
normalization one has \ldots'' \hfill\anchor{ATS~IV}{3185--3187}
\end{quote}

prepared in the introduction by

\begin{quote}\small
``Each arithmeticoid $y_0' = \phi(y_0)$, $y_0$ provides an isomorphic copy of our
number field and hence one has chosen the same global normalization for both the
copies, the [sic] it does not matter which arithmeticoid is used for global
computations.'' \hfill\anchor{ATS~IV}{740--743}
\end{quote}

On the other hand, immediately below the theorem, the same text states:

\begin{quote}\small
``as arithmeticoids $y_0' \neq y_0$ as they differ by Frobenius at each prime
$v \in V_L$ and this, together with choice of concurrent global normalization for
the copy of the (same) number field they provide \emph{precludes direct comparison
of various local quantities}.'' \hfill\anchor{ATS~IV}{3207--3211}
\end{quote}

Both passages are Joshi's. They motivate a question about the identification
of number-field element scales, but do not place the entire E2 equality on
that identification. The indexed sets, algebraic hulls, measures and weights
require their own bindings.

\begin{question}\label{q:main}
Does the phrase ``same global normalization'', as used in ATS~IV Theorem~6.10.1,
\emph{assert} an isometry between the two valuation scales, or does it
\emph{establish} one by definition?
\end{question}

The two answers have very different consequences. If the phrase asserts an
isometry, then something is being claimed that a reader is entitled to see proved,
and the program's working name for the missing statement --- a \emph{global
normalization contract lemma} --- would be justified. If the phrase establishes
the element-scale identity definitionally, then no separate lemma is needed
for that limited identity. What carries the indexed hull and volume equality
is a further question.

This paper answers Question~\ref{q:main} in the second sense --- conditionally on
a reading that the sources leave open, as Section~\ref{sec:trilemma} sets out ---
and then pursues the follow-up question, because the follow-up is where the
substance turned out to be.

\subsection{Sources, labels, and citation discipline}
\label{sec:sources}

Four documents are used as primary text. Throughout we use the short labels in
the left column; every source citation in this paper carries such a label,
because the four extracts have independent line numbering and a bare line number
would be ambiguous. Table~\ref{tab:sources} summarizes the primary source corpus and assigned roles.

\begin{table}[htbp]
\centering
\caption{Primary source corpus and assigned labels in this assessment.}
\label{tab:sources}
\small
\begin{tabular}{@{}lll@{}}
\toprule
Label & Source & Role in this paper \\
\midrule
ATS~II & arXiv:2303.01662v3~\cite{JoshiATSII}
  & behavior of valuations along the Frobenius fiber \\
ATS~II(1/2) & arXiv:2305.10398v12~\cite{JoshiATSIIhalf}
  & definition of $\arith(L)^{\mathrm{nor}}$; period mapping; heights \\
ATS~III & arXiv:2401.13508~\cite{JoshiATSIII}
  & normalization convention; hyperplane; volume theory \\
ATS~IV & arXiv:2403.10430~\cite{JoshiATSIV}
  & Theorem~6.10.1 and its surrounding apparatus \\
\bottomrule
\end{tabular}
\end{table}

In Joshi's own bibliographies these appear as ``[Joshi, 2023b]'' (ATS~II),
``[Joshi, 2023a]'' (ATS~II(1/2)) and ``[Joshi, 2024c]'' (ATS~III); the assignment
of the first two was verified against the bibliography of ATS~IV
\anchor{ATS~IV}{3807--3810}, the third is the identification used consistently in
the working notes underlying this paper. Where we reproduce a passage, the wording
is verbatim from the extract; mathematical symbols have been typeset, and
ellipses are marked. Two OCR-level readings were made explicitly: the string
``ptvh-power'' in ATS~II(1/2) Remark~5.9.2 is read as ``$p_v$-th power'', and the
subscript/superscript placements in the displayed normalization equations were
checked against their surrounding prose.

\subsection{What follows}

Section~\ref{sec:equivariance} sharpens the question: the naive form ``is $F$
invariant?'' misses the structure, because both the evaluated object and the scale
in which it is evaluated move. Section~\ref{sec:definition} establishes the meaning
of the phrase under the place-wise reading of the definition of
$\arith(L)^{\mathrm{nor}}$; its
Section~\ref{sec:trilemma} then states the trilemma that the definition leaves
open, and Section~\ref{sec:rigidity} proves the lemma that constrains it.
Section~\ref{sec:D1} shows by an elementary derivation, under two assumptions
stated with it, that the product-formula hyperplane really does move, which
excludes --- under those assumptions --- the one reading under which the phrase
would have had non-trivial content.
Section~\ref{sec:terms} classifies the terms of Theorem~6.10.1 and locates the
candidate carrier. Section~\ref{sec:measure} separates local conventions from
the open cross-frame measure and hull obligations. Section~\ref{sec:fneu} states the open question at
the margin --- the one that the trilemma grew out of.
Section~\ref{sec:assessment} draws the consequences and states the limitations.
Verdicts that depend on the branch are marked as such where they are stated; the
summary of which verdict holds under which branch is Table~\ref{tab:branches}.

% =============================================================================
\section{Sharpening the question: equivariance, not invariance}
\label{sec:equivariance}
% =============================================================================

\subsection{Both sides move}

A first attempt to formalize Question~\ref{q:main} would ask: is the quantity $F$
invariant under a change of arithmeticoid? That formulation is wrong, and seeing
why is the first substantive step.

In the absolute-value equality (with the indexing detailed in E2 below)
\[
 \bigl|\log\Vol(\operatorname{hull}(\Thet^{\varphi(y_0)}))\bigr|
 =\bigl|\log\Vol(\operatorname{hull}(\Thet^{y_0}))\bigr|
\]
\emph{both} arguments move: the evaluated object and the scale in which the
evaluation happens. The correct form of the statement is therefore not invariance
but an indexed, frame-dependent comparison. An isometric transport would be
 a possible sufficient route; the absolute-value equality alone does not imply it. This matters in
both directions. A statement of this type is not automatically trivial --- being
an isometry is a substantive property --- but it is also not automatically empty
when it holds, unlike a mere renaming. Whether the phrase ``identically
normalized'' asserts such an isometry or produces it by an identification is
exactly Question~\ref{q:main}.

The relevant structural facts are stated by Joshi himself. The schemes are
isomorphic:

\begin{quote}\small
``the schemes underlying these holomorphoids are isomorphic to $X/L$''
\hfill\anchor{ATS~IV}{426--427}
\end{quote}

while the analytic structures, and with them the valuation scales, are not:

\begin{quote}\small
``arithmetic (and geometry) provided by the holomorphoids $\hol(X/L)_y$ and
$\hol(X/L)_{y'}$ \ldots\ occur at intrinsically different valuation scales''
\hfill\anchor{ATS~IV}{450--453}
\end{quote}

and this difference is not an inconvenience of the construction but its purpose:

\begin{quote}\small
``It is precisely the existence of such intrinsic differences in valuation scales
which the averaging techniques \ldots\ exploit''
\hfill\anchor{ATS~IV}{457--459}
\end{quote}

\subsection{Two precision levels of the question}

Two versions of the question can be stated, of unequal maturity.

\paragraph{(L-weak) The invariance class.}
Let $N$ be a fixed choice of identical global normalization for the copies of $L$
provided by $y_0$ and $\varphi(y_0)$. Characterize the class
\[
  G \;:=\; \bigl\{\, F \;:\; F(y_0) = F(\varphi(y_0)) \text{ is \emph{provable}
  under } N \,\bigr\},
\]
where $F$ is evaluated covariantly (object and scale transported together) and
``provable'' explicitly \emph{excludes} that the equality follows from the
definition of $N$ or from a notational convention.

A proof that $F \in G$ would have the shape of a chain of equalities carrying
$F(\varphi(y_0))$ into $F(y_0)$ using only steps that either run over data
demonstrably fixed by $\varphi$, or exploit a demonstrated compensation of two
changes. The text supplies a formal model for the first kind: the proof of ATS~IV
Proposition~4.4.4 reduces $\log(q_M) = \log(q_{L_{\mathrm{mod}}})$ to the degree
formula $\sum_{w\mid v} e_{w\mid v} f_{w\mid v} = [M : L_{\mathrm{mod}}]$
\anchor{ATS~IV}{1866--1928}. That is a genuine proof by classical means --- but it
establishes invariance under \emph{field extension}, not under $\varphi$, and is
therefore a template here, not an answer.

Several source passages delimit the proposed class. ATS~IV reports that
Tate-parameter absolute values rise under Frobenius
\anchor{ATS~IV}{567--569}, that the normalized Tate-idele degree is
non-constant on $\Ynum$ \anchor{ATS~IV}{2055--2075}, and that height
is written on $\overline L^*\times\Ynum$
\anchor{ATS~IV}{550--559}. ATS~III reports movement of the
product-formula hyperplane \anchor{ATS~III}{518--522}.
These are source-reported examples and indications, not four independently
proved failures of Frobenius invariance. A displayed domain alone proves
neither non-constancy nor variation along a particular Frobenius orbit.
The height statement and F-NEU-1 retain the separate domain and proof
obligations of Section~\ref{sec:fneu}; D-1 retains (F1)--(F2).

\paragraph{(L-strong) Order compatibility after summation.}
Give a criterion under which global invariance yields \emph{order compatibility}
between two families of local contributions, although place-by-place
comparability is excluded. The side condition is not invented; it is Joshi's:

\begin{quote}\small
``one cannot claim that $z_p \le z_p'$ holds for all primes $p$ without violating
the middle equality of (1.5.1)'' \hfill\anchor{ATS~IV}{817--819}
\end{quote}

This second version is the less developed of the two: it names the structure but
offers no candidate class of admissible local distributions. We record that as the
largest substantive arrear of the formulation phase; it is not resolved in this
paper.

\subsection{A separating criterion}

The classification used in Sections~\ref{sec:terms} and~\ref{sec:measure} rests on
a criterion that we state as a hypothesis, not as a theorem.

\begin{criterionK}
A quantity is the more likely to be invariant in the load-bearing sense, the more
completely it factors over data that $\varphi$ demonstrably fixes (the underlying
scheme, field degrees, residue field cardinalities), and the less so, the more
directly it reads the \emph{valuation scale} $|\cdot|_{K_{y_v}}$ of the
arithmeticoid.
\end{criterionK}

The criterion's discriminating power can be read off the theorem itself, because
Theorem~6.10.1 contains both types side by side: its first equality concerns
$q_\ell$, whose defining formula runs over $\ordd_w(q_w)$ and $\log w$, while its
middle equality concerns $\log \Vol$, which through
$\LogVol(\pi \mathcal{O}) = \log|\pi|$ reads the scale directly. Separating the
two was the first result of the formulation phase; carrying the separation through
is the business of Section~\ref{sec:terms}. Table~\ref{tab:invariance_classes}
summarizes the resulting taxonomy of invariance classes under an arithmeticoid shift.

\begin{table}[htbp]
\centering
\caption{Taxonomy of invariance classes and the separating Criterion~K under an arithmeticoid shift $\varphi$.}
\label{tab:invariance_classes}
\small
\setlength{\tabcolsep}{3.5pt}
\begin{tabular}{@{}P{0.07\textwidth}P{0.24\textwidth}P{0.24\textwidth}P{0.23\textwidth}P{0.14\textwidth}@{}}
\toprule
Class & Invariance Character & Factorization Basis & Exemplars in ATS~IV & Gap Role \\
\midrule
\textbf{T} & Trivial / definitional & Factors over $\varphi$-fixed scheme \& degree data ($[M:\mathbb{Q}]$, $e, f$) & $\ordd_w(q_w)$, $\log w$, $q_\ell$ (\S\,4.4) & Harmless \\
\textbf{B} & Bound-invariance (global) & Preserved under global summation / product formula & $\sum_v \log |x|_v = 0$, $A' = CA$ & Constraining \\
\textbf{F} & Fiber-equivariance & Transforms equivariantly along Frobenius fiber & Period map $\operatorname{per}_{\arith}$, $\Thet$ & Structural \\
\textbf{N} & Non-invariance (untilt) & Directly evaluates untilt valuation scale $|\cdot|_{K_{y_v}}$ & Tate idele $\log(TI)$, Kummer map & Obstructing \\
\bottomrule
\end{tabular}
\end{table}

% =============================================================================
\section{The definition under the place-wise reading:
\texorpdfstring{$\arith(L)^{\mathrm{nor}}$}{arith(L)nor}
is unique per point}
\label{sec:definition}
% =============================================================================

\subsection{Three candidate readings (not exhaustive)}

Before consulting the definition, the phrase ``same global normalization'' admits
at least three readings. It is useful to state these three because two can be
excluded on grounds internal to the text; the list is diagnostic, not a claim of
logical exhaustiveness.

\begin{description}
\item[N1.] Both arithmeticoids are normalized in the sense of the standing
  convention, i.e.,\ the product formula holds in each.
\item[N2.] Both have the \emph{same} product-formula hyperplane,
  $H_{y_0} = H_{\varphi(y_0)}$.
\item[N3.] A coordinated choice is made within the \emph{residual freedom} that
  the normalization condition leaves open --- the thought being that
  $\prod_v |x|_{K_v} = 1$ does not fix the individual $|\cdot|_{K_v}$.
\end{description}

For example, a further mechanism could supplement the product-formula condition
with a global selection rule that fixes the residual scalar $c$, without returning
to the place-wise Artin--Whaples binding of branch~(a)/N1$'$ and without allowing
the place-selective freedom of branch~(c)/N3. These three readings therefore organize the
source passages at hand; they do not partition every possible selection rule.

Reading N1 cannot be what is meant on its own, because the convention is already
universal in the series:

\begin{quote}\small
``I will also assume throughout this paper that all arithmeticoids $\arith(L)_y$
are normalized i.e.\ $\arith(L)_y = \arith(L)^{\mathrm{nor}}_y$ \ldots\ i.e.\ the
valuations on each component of $y$ are normalized so that the Artin--Whaples
product formula (\cite{ArtinWhaples1945}) holds for each $x \in L^*$:
$\prod_{v \in V_L} |x|_{K_v} = 1$'' \hfill\anchor{ATS~III}{858--863}
\end{quote}

Under N1, the sentence in Theorem~6.10.1 would add nothing that does not hold
anyway. Reading N2 stands in visible tension with the text's own emphasis on the
mobility of the hyperplane, to which we return in Section~\ref{sec:D1}. That
leaves N3 as the reading that would carry content --- and N3 is precisely the one
whose definition is not stated in ATS~III or ATS~IV, but in the earlier paper.

\subsection{What the definition actually says}

The definition of the normalized arithmeticoid is given in ATS~II(1/2)
\S\,5.3. The relevant passages are:

\begin{quote}\small
``The fundamental global constraint \ldots\ is the validity of the product formula
for $L$ \ldots\ for the \emph{standard choice of normalizations} for valuations
$(L_v, |\cdot|_v)$ (given in [Artin and Whaples, 1945]) \ldots\ (5.3.1)
$\prod |x|_v = 1$'' \hfill\anchor{ATS~II(1/2)}{1840--1848}
\end{quote}

\begin{quote}\small
``Now suppose one is given a $y = (\psi_v) \in \Ynum$ \ldots\ Since the two
valuations $|\cdot|_v$ and $|\cdot|_{K_{\psi_v}}$ on $L_v$ are \emph{equivalent,
there exist, for each $v \in V_L$, real numbers $\lambda_v \in \mathbb{R}$} such
that (5.3.3) $(L_v, |\cdot|_v) = (L_v, |\cdot|^{\lambda_v}_{K_{\psi_v}})$.''
\hfill\anchor{ATS~II(1/2)}{1852--1858}
\end{quote}

\begin{quote}\small
``(5.3.5) $\lambda_y = (\lambda_v)$ \ldots\ \emph{is the normalization coordinate
of $y$}. When valuations induced on $\{L_v\}$ by the local data of $\arith(L)_y$
are normalized so that (5.3.4) holds, I will write this as
$\arith(L)^{\mathrm{nor}}_y$'' \hfill\anchor{ATS~II(1/2)}{1867--1873}
\end{quote}

\begin{quote}\small
``Because the absolute values $(|\cdot|_{K_{\psi_v}})$ depend on $y \in \Ynum$,
one has (5.3.7) $\lambda_{y'} \neq \lambda_y$ in general and hence one says that
the \emph{product formula normalization of $y$ moves with $y$} $\in \Ynum$.
Especially, one sees that there is \emph{no uniform choice} of the normalization
coordinate $\lambda_y$ for all $y \in \Ynum$.''
\hfill\anchor{ATS~II(1/2)}{1884--1891}
\end{quote}

Two further pieces of context belong here. Each point $y$ carries a
\emph{canonical} valuation, obtained via the Fargues--Fontaine curve~\cite{FarguesFontaine2018},
and alongside it the \emph{$L_v$-standard normalized}
valuation with $|p|_{K_v} = 1/p^{f_v}$ \anchor{ATS~II(1/2)}{802--840}. And the
impossibility of normalizing everything at once is stated as a statement about the
family:

\begin{quote}\small
``Regardless of the choice of a normalization \ldots\ one cannot normalize
valuations of all untilts of $\bar{L}_v$ \emph{simultaneously}. More precisely the
function $y \mapsto |p|_{K_y}$ is not a constant function.''
\hfill\anchor{ATS~II(1/2)}{802--840, Rem.~2.7.1}
\end{quote}

\subsection{Three consequences}

\paragraph{(i) Read as a place-wise binding, (5.3.3) determines the normalization
uniquely for each fixed $y$.}
The exponent $\lambda_v$ is fixed by (5.3.3) as the conversion exponent between
the canonical $K_{\psi_v}$-valuation and the Artin--Whaples standard valuation.
This is the classical uniqueness statement for equivalent absolute values: if
$|\cdot|_1$ is non-trivial and $|\cdot|_2 = |\cdot|_1^{\,a}$ for some $a > 0$,
then $a$ is unique, being computable as $\log|y|_2/\log|y|_1$ for any $y$ with
$|y|_1 > 1$ \cite[Prop.~7.8, p.~108]{Milne2020}. Read this way there is \emph{no
residual freedom} within which anything could be ``coordinated''. That reading is
natural --- and it is the one this section's verdict rests on --- but it is not
the only one the sources admit; Section~\ref{sec:trilemma} states the alternative.

\paragraph{(ii) On elements of the number field the normalized valuations agree
with the standard ones, for every $y$.}
That is the content of (5.3.3)/(5.3.4). What moves with $y$ is the
\emph{coordinate} $\lambda_y$ --- the position of the standard scale relative to
the canonical untilt scale --- not the value $|x|$ for $x \in L$.

\paragraph{(iii) The non-simultaneity is a statement about the family.}
Remark~2.7.1 denies a $y$-uniform normalization across $\Ynum$; it is not an
obstruction to the pointwise normalization of any individual $y$. Distinguishing
the family statement from the point statement is what makes (i) and the quoted
``no uniform choice'' compatible rather than contradictory.

\subsection{Decision of the trichotomy, under the place-wise reading}
\label{sec:trichotomy}

\begin{itemize}
\item \textbf{N3 is vacuous} --- if (5.3.3) is read as a place-wise binding.
  Section \S\,5.3 then leaves no residual freedom: the freedom conjectured in the
  formulation phase --- that $\prod_v |x|_v = 1$ does not pin the individual
  absolute values --- does not exist, because it is not \emph{some}
  product-formula-compatible family that is chosen, but \emph{the} one bound to
  the Artin--Whaples standard by (5.3.3). Section~\ref{sec:trilemma} shows that
  this qualification is not idle: elsewhere in the series a residual freedom of
  exactly the N3 kind is \emph{exercised}.
\item \textbf{N2 is excluded under the assumptions (F1)--(F2)} of the derivation
  of Section~\ref{sec:D1}; without them it is textually disfavored but open. The
  exclusion is branch-independent --- it uses only the movement of $\lambda_y$
  along a Frobenius fiber, which every branch retains --- but it is not
  assumption-free: (F1)--(F2) place the standard point of ATS~IV on such a fiber
  and supply the coordinate transformation at two finite places.
\item \textbf{N1$'$ holds under that reading and is then the meaning of the
  phrase.} By N1$'$ we mean a sharpened form of N1: both copies are taken in their
  (individually unique) normalized form, whereby elements of the number field
  receive the same standard absolute values in both copies. This is the only thing
  that ``Choose the same global normalization for the number field provided by the
  arithmeticoids $y_0'$ and $y_0$'' can mean, given a place-wise reading of
  \S\,5.3, and it is exactly what Joshi's own justification says: the
  arithmeticoids ``provide an isomorphic copy \ldots\ hence one has chosen the same
  global normalization''.
\end{itemize}

The wording as a ``choice'' is, in this light, mildly misleading: there would be
nothing to choose. It would be the invocation of the standing convention of
ATS~III, under which all arithmeticoids are normalized. Question~\ref{q:main}
would thereby be answered in its second alternative --- the phrase establishes the
identity by definition and does not assert an isometry --- and on the level of
number-field elements a ``missing contract lemma'' of the conjectured form would
not exist.

We write this in the conditional deliberately. The verdict is sound \emph{given}
the place-wise reading of (5.3.3), and that reading is the natural one for the
wording of \S\,5.3 taken by itself. It is not, however, the reading under which
the normalization is used in the rest of the series. The next subsection sets out
the resulting alternative, and every later verdict of this paper is marked for the
branches under which it holds.

Independently of the branch, the definition identifies \emph{nothing} on the
untilt level, where the canonical valuations of neighboring points on a Frobenius
fiber demonstrably differ already on $\mathbb{Q}_p$. Which of the quantities in
Theorem~6.10.1 live on which level is the subject of Section~\ref{sec:terms}.

% -----------------------------------------------------------------------------
\subsection{What the definition does not state: a trilemma for
\texorpdfstring{$\lambda_y$}{lambda\_y}}
\label{sec:trilemma}
% -----------------------------------------------------------------------------

The preceding subsection assumed an answer to a question the sources do not
address: \emph{which rule selects the normalization coordinate $\lambda_y$ when
the arithmeticoid moves?} We searched the passages listed in
Section~\ref{sec:sources} for such a rule and did not find one. What we did find
is that the two candidate answers are both used --- in different places.

\paragraph{The place-wise reading is what \S\,5.3 says.}
Equation (5.3.3) fixes $\lambda_v$ as the exponent converting the canonical
valuation into the Artin--Whaples standard one, and (5.3.5) speaks of
\emph{the} normalization coordinate. On that reading (5.3.4) is a rewriting of
(5.3.1) rather than an additional constraint --- the text itself says the product
formula ``reads'' that way.

\paragraph{A place-selective freedom is what the later papers exercise.}
ATS~III states that the scaling at the semi-stable places compels a
re-normalization \emph{at the other places}:

\begin{quote}\small
``the non-trivial scaling relationship established by Theorem~4.2.2.1(3) at primes
$w \in V_{\mathrm{odd},ss}$ \emph{forces that suitable (re)normalization of
valuations must be introduced at primes} $w \in V_{L'} - V_{\mathrm{odd},ss}$ so
that the product formula to continue to hold for each normalized arithmeticoid
$\arith(L)^{\mathrm{nor}}_{y'_j}$.'' \hfill\anchor{ATS~III}{1600--1609}
\end{quote}

The same mechanism appears as a structural property in Theorem~4.6.1(4)
\anchor{ATS~III}{1786--1791}, and ATS~IV explains the movement of heights by it:
the local contributions are ``\emph{bound together} \ldots\ by means of the
(global) product formula'', so that the global Frobenius ``changes global heights
in a subtle and complicated way in order to ensure that the product formula
holds'' \anchor{ATS~IV}{585--589}. An instruction to introduce a re-normalization
at other places presupposes that at those places something is still free to be
chosen --- which is precisely the freedom that a place-wise reading of (5.3.3)
denies.

\paragraph{Three diagnostic branches (not exhaustive).}
Since no rule is stated, three source-motivated possibilities must be kept apart:

\begin{description}
\item[(a) Standard-bound.] $\lambda_y$ \emph{is} the tuple determined place by
  place by (5.3.3). Then $\arith(L)^{\mathrm{nor}}_y$ is well defined, and
  Section~\ref{sec:trichotomy} applies as written --- but there is nothing to
  ``introduce'' at the other places, and the instruction quoted above has no
  object.
\item[(b) Product-formula-bound.] $\lambda_y$ is any tuple for which the product
  formula holds. Then it is well defined only up to a scalar, by Lemma~1 below,
  and the freedom is \emph{not} place-selective.
\item[(c) Place-selective.] $\lambda_y$ is selected by some further rule that
  leaves genuine place-by-place freedom. This is what the passages just quoted
  use --- but the rule itself is not given in the material we read, and it is not
  reconcilable with reading (5.3.3) as a place-wise binding.
\end{description}

Whether (a), (b) or (c) --- or a further global selection rule such as the one
noted above --- is intended cannot, as far as we can see, be decided from the
text. It is the natural first question of any correspondence about this material,
and we formulate it as such in Section~\ref{sec:fneu}.

% -----------------------------------------------------------------------------
\subsection{What the product formula alone can supply: a rigidity lemma}
\label{sec:rigidity}
% -----------------------------------------------------------------------------

Branch (b) can be settled. The product formula constrains the multiplier tuple far
more strongly than one might expect: it determines it completely, up to one global
scalar.

\paragraph{Conversion multipliers used by the lemma.}
Let $\lambda^{\mathrm{std}}_{y,v}>0$ satisfy
$|x|_{K_{y_v}}^{\lambda^{\mathrm{std}}_{y,v}}=|x|^{\mathrm{std}}_v$ on $L_v$,
and write $\lambda_{y,v}=\mu_v\lambda^{\mathrm{std}}_{y,v}$.
Then $|x|_{K_{y_v}}^{\lambda_{y,v}}=(|x|^{\mathrm{std}}_v)^{\mu_v}$.
The following lemma constrains $\mu$, not directly $\lambda_y$.
Branch~(a) fixes $\mu_v=1$; branch~(b) restricts positive conversion exponents
to $\mu_v=c>0$. It makes no corresponding assertion about Haar measures.

\begin{lemmaR}[Rigidity of the product-formula multiplier]
\label{lem:rigidity}
Let $L$ be a number field, and for each place $v$ let $|\cdot|^{\mathrm{std}}_v$
denote the Artin--Whaples standard absolute value, so that
$\sum_v \log|x|^{\mathrm{std}}_v = 0$ for all $x \in L^*$. Let
$(\mu_v)_{v \in V_L}$ be real numbers with
\[
  \sum_{v \in V_L} \mu_v \cdot \log|x|^{\mathrm{std}}_v \;=\; 0
  \qquad \text{for all } x \in L^*.
\]
Then all $\mu_v$ are equal.
\end{lemmaR}

\noindent\emph{Proof sketch.}
For $L = \mathbb{Q}$ the statement is immediate and needs nothing beyond the
standard absolute values. Substituting $x = p$ for a prime $p$ gives
$\log|p|_p = -\log p$, $\log|p|_\infty = +\log p$ and $\log|p|_q = 0$ for primes
$q \neq p$, so the hypothesis reads $(\mu_\infty - \mu_p)\log p = 0$, i.e.,\
$\mu_p = \mu_\infty$ for every prime $p$; $x = -1$ adds no condition, and
$\mathbb{Q}^*$ is generated by $-1$ and the primes.

For a general number field the same two steps are used, with two classical inputs.
First, for a unit $x \in \mathcal{O}_L^*$ all finite terms vanish, and by
Dirichlet's unit theorem the vectors $(\log|x|_v)_{v \text{ arch}}$ span the
trace-zero hyperplane of the archimedean coordinates; hence $\mu_v$ is constant
$= c$ over the archimedean places. Second, for a finite place $v$ let $h$ be the
class number of $L$ and $\pi$ a generator of the principal ideal
$\mathfrak{p}_v^{\,h}$; substituting $x = \pi$ gives
$-\mu_v \cdot h\log N(\mathfrak{p}_v) + c \cdot h\log N(\mathfrak{p}_v) = 0$, i.e.,\
$\mu_v = c$. \hfill$\square$

\paragraph{Status and provenance.}
The lemma is elementary and certainly not new in substance; we state it because
the constraint it expresses is used implicitly in the material and, as far as we
could determine, is nowhere cited there. The $\mathbb{Q}$-case above is complete
as it stands. The general case uses Dirichlet's unit theorem and the finiteness of
the class number, both classical; we give no section number for either, because we
did not verify one in an edition we hold. We deliberately do \emph{not} attribute
the lemma to~\cite{ArtinWhaples1945}: the product-formula characterization proved
there concerns which fields admit such a system of absolute values, and we could
not obtain the full text to check whether a uniqueness statement about the
multipliers is included. Anyone citing this lemma should either prove it, as here,
or verify a source directly.

\paragraph{Consequence for branch (b), stated with care.}
By Lemma~\ref{lem:rigidity} the admissible tuples in branch (b) are exactly
$\mu \equiv c$ with $c > 0$. Two things follow, and they pull in opposite
directions.

First, the freedom that branch (b) grants is \emph{uniform}, never selective. A
global factor multiplies every local contribution alike; it cannot raise the
contribution at one place and lower it at another. Wherever the material describes
a re-normalization at the non-semi-stable places as \emph{compensating} a scaling
at the semi-stable ones, that description asks for more than the product formula
by itself can deliver. This is a statement about what one constraint implies --- it
is not a claim that Corollary~4.2.2.4 is false, and not a claim that no rule
achieving the described effect exists. It says that such a rule would have to come
from somewhere other than the product formula, and it thereby sharpens what a
resolution of the trilemma would have to supply.

Second, the scalar $c$ is not fixed by the constraint. For a projective object
such as the hyperplane $H_y$ this is harmless --- $c$ cancels, since coefficient
vectors define the same hyperplane exactly when they are proportional --- but for
a quantity that is linear in $\lambda$, such as a height, it is not: the value
changes by the factor $c$. Branch (b) therefore leaves the scale of such
quantities undetermined unless a further convention is added.

\begin{table}[htbp]
\centering
\caption{Which verdict of this paper holds under which branch of the trilemma of
Section~\ref{sec:trilemma}. ``---'' marks a verdict that the branch does not
support as stated. Branch-independence is not the same as
assumption-freedom: the first row holds in every branch \emph{under the
assumptions} (F1)--(F2) of Section~\ref{sec:D1}, and the fifth row identifies the
carrier of the middle equality without establishing the transport to the two
indexed hull evaluations displayed in (E2) (Section~\ref{sec:residues}).}
\label{tab:branches}
\small
\begin{tabular}{@{}P{0.22\textwidth}P{0.15\textwidth}P{0.15\textwidth}P{0.17\textwidth}P{0.10\textwidth}@{}}
\toprule
Verdict & (a) standard-bound & (b) PF-bound & (c) place-selective & \shortstack{Evidence\\level} \\
\midrule
N2 excluded (D-1, under (F1)--(F2)) & holds & holds & holds & \textbf{O} \\[3pt]
N3 vacuous & holds & partly: freedom exists but is not selective & --- & \textbf{O} \\[3pt]
Scale identification definitional (N1$'$) & holds & up to the scalar $c$ & --- & \textbf{D} \\[3pt]
E1 ($q_\ell$) invariant & holds & holds & holds (reads no scale) & \textbf{D} \\[3pt]
E2 candidate carrier ($\varphi$-stability of $\Thet$; indexed transport untraced)
  & holds & holds & holds & \textbf{S} \\[3pt]
E2 measure compatibility & open & open & open & \textbf{O} \\[3pt]
Nothing identified on the untilt level & holds & holds & holds & \textbf{D} \\
\bottomrule
\end{tabular}
\par\smallskip
{\footnotesize \textbf{D}: definitional;\quad \textbf{S}: source-asserted;\quad
\textbf{O}: open or conditional.\par}
\end{table}

% =============================================================================
\section{The period mapping really moves: the conditional exclusion of N2}
\label{sec:D1}
% =============================================================================

\subsection{The mobility of the hyperplane in the sources}

Each normalized arithmeticoid provides a degree homomorphism and, through the
logarithm of the product formula, a hyperplane:

\begin{quote}\small
``(2) Each normalized arithmeticoid $\arith(L)_y$ provides a degree homomorphism
\ldots\ $(x_v) \mapsto \sum_{v \in V_L} \log(|x_v|_{K_{y_v}}) = \deg_{y_j}((x_v)_{v \in V_L})$.
(3) As a consequence, each normalized $\arith(L)^{\mathrm{nor}}_y$ provides a
hyperplane $H_{y_j} \subset V_L$ given by the logarithm of product formula equation
(2.1.1)'' \hfill\anchor{ATS~III}{1762--1783}
\end{quote}

The hyperplane is not fixed:

\begin{quote}\small
``the above arithmetic period mapping is a highly non-trivial global arithmetic
period mapping in which an arithmetic avatar $\arith(L)_y$ provides a hyperplane
$H_y$ (as above) \emph{which moves under symmetries (such as the actions of the
absolute Galois group, global Frobenius and $L^*$)}''
\hfill\anchor{ATS~III}{518--522}
\end{quote}

and, correspondingly, ``the global arithmetic period mapping is non-constant''
\anchor{ATS~III}{869--870}. The reason given is the interplay of local freedom and
global constraint:

\begin{quote}\small
``the valuations on untilts parameterized by a Fargues--Fontaine curve at some
arbitrary prime \ldots\ cannot be all simultaneously normalized and on the other
hand, local changes of valuations are globally constrained by the validity of the
product formula'' \hfill\anchor{ATS~III}{515--518}
\end{quote}

In the earlier paper the hyperplane is given \emph{explicitly with the
normalization coordinates as coefficients}:

\begin{quote}\small
``$H_y = \{ z = (z_v) \in \bigoplus_v \mathbb{R} : \sum_v \lambda_v z_v = 0 \}$''
\hfill\anchor{ATS~II(1/2)}{2199--2206}
\end{quote}

together with the statement that ``one has a natural and \emph{non-constant}
mapping $y \mapsto H_y$ \ldots\ called the period mapping''
\anchor{ATS~II(1/2)}{2227--2233}, whose non-constancy is proved as being
``immediate from the fact that \emph{normalizations of arithmeticoids do not carry
over} because locally at each prime $v$, valuations \emph{cannot be simultaneously
normalized}'' \anchor{ATS~II(1/2)}{2238--2243}. Remark~5.10.2(5) of the same paper
notes that the impossibility objection of Scholze--Stix~\cite{ScholzeStix2018}
addresses exactly this phenomenon; ATS~III likewise records that Scholze--Stix
dispute the mobility of $H_y$ \anchor{ATS~III}{522--525}. We report this
disagreement and take no side in it.

\subsection{The behavior along the Frobenius fiber}

On the fiber $\{y_n\}_{n \in \mathbb{Z}}$ over the canonical point --- the
deck-transformation, i.e.,\ Frobenius, orbit --- the sources state:

\begin{itemize}
\item the Kummer embeddings $\mathbb{Z}_p(1) \to B$ attached to $y_n$ and
  $y_{n-1}$ ``differ as $\mathbb{Z}_p$-submodules by a $p$-multiple''
  \anchor{ATS~II}{2119--2120};
\item ``One has $|p|_{K_{y_{n-1}}} = |p|^{1/p}_{K_{y_n}}$''
  \anchor{ATS~II}{2122};
\item ``the valuations of $p$, \emph{and hence of the Tate parameters} of the
  holomorphoids \ldots\ \emph{grow along the fiber} $\{y_n\}$''
  \anchor{ATS~II}{2127--2131};
\item ``each $y_n$ comes with a canonical valuation on its residue field
  $K_{y_n}$ \ldots\ but the valuations provided by $y_{n-1}$ and $y_n$ \emph{are
  not the same on $\mathbb{Q}_p$}'' \anchor{ATS~II}{2184--2188}.
\end{itemize}

and, for the action at all places at once, that the global Frobenius relates the
local multiplicative structures of $L$ and $L^{(1)}$ ``through the \emph{$p_v$-th
power mapping for each $v \in V_L$}'' \anchor{ATS~II(1/2)}{2149--2153}.

\subsection{The derivation}

\begin{derivationD}
Assume \textbf{(F1)} that the standard point $y_0$ used in ATS~IV \S\,4.1 lies in a
Frobenius fiber to which ATS~II Theorem~10.20.1(3) applies, and \textbf{(F2)} that
ATS~II(1/2) Remark~5.9.2 induces $\lambda_{\varphi(y),v} = p_v \lambda_{y,v}$ at two
finite places of distinct residue characteristic. Then
$H_{\varphi(y)} \neq H_y$: the product-formula hyperplane is not preserved by the
global Frobenius.
\end{derivationD}

\noindent\emph{Derivation.}
\begin{enumerate}
\item $H_y$ has coefficient vector $\lambda_y$
  \anchor{ATS~II(1/2)}{2199--2206}.
\item Along the canonical fiber one has
  $|p|_{K_{y_{n-1}}} = |p|^{1/p}_{K_{y_n}}$ \anchor{ATS~II}{2122}. With the
  standard condition $|p|^{\lambda_n}_{K_{y_n}} = 1/p$ (case $L = \mathbb{Q}$,
  $v = p$) this yields $\lambda_{n-1} = p \cdot \lambda_n$.
\item At \emph{every} place $v$, $\varphi$ acts as a $p_v$-th power step
  \anchor{ATS~II(1/2)}{2149--2153}, i.e.,\ $\lambda_v \mapsto p_v \cdot \lambda_v$.
  The step from the quoted statement --- which speaks of the local
  \emph{multiplicative structures} --- to this statement about the
  \emph{normalization coordinates} is read off the quotation by analogy with
  step 2; it is carried out as a computation only in the case $L = \mathbb{Q}$,
  $v = p$ there. This is assumption (F2).
\item Two coefficient vectors define the same hyperplane only if they are
  proportional by a \emph{single} scalar $c$. But $(p_v \lambda_v) = c(\lambda_v)$
  would force $p_v = c$ for all $v$, which is impossible because the residue
  characteristics vary. Hence $H_{\varphi(y)} \neq H_y$. \hfill$\square$
\end{enumerate}

\paragraph{Status of this derivation.}
D-1 is \emph{our own} derivation --- elementary algebra plus a proportionality
argument --- assembled from statements quoted verbatim above. It is \emph{not} a
theorem of the source, and it is not presented as one. It is stated here because
it settles reading N2, which the formulation phase had been able to classify only
as ``documented-unlikely''.

\paragraph{Two caveats, carried explicitly.}
\begin{enumerate}
\item[(i)] Step 2 computes in the canonical fiber over the canonical point. That
  the standard points $y_0$ fixed in ATS~IV \S\,4.1 lie on such a fiber is
  consistent with the statement that ``there are countably infinite possible
  choices for the standard point $y_0$ \ldots\ Any one of the countably many
  choices for $y_0$ are equally valid choices'' \anchor{ATS~III}{1723--1729}, but
  it was not separately established here.
\item[(ii)] The behavior at archimedean places is asserted in Remark~5.9.2 by the
  phrase ``for each $v \in V_L$'', but was not separately checked. For D-1 this is
  not needed: two finite places with distinct residue characteristics already
  suffice for step 4.
\end{enumerate}

Under these two caveats --- which are the assumptions (F1)--(F2) of the statement
above --- N2 is excluded rather than merely improbable, and the
fixed-point reservation that the formulation phase had to carry --- the
possibility that $y_0$ might happen to be a $\varphi$-fixed point --- is relieved:
along canonical fibers the normalization coordinate demonstrably moves, and a
fixed standard point would have to lie outside such fibers, for which nothing in
the material speaks.

% =============================================================================
\section{Term classification of Theorem~6.10.1: three kinds of statement}
\label{sec:terms}
% =============================================================================

\subsection{The chain}

In simplified notation, Theorem~6.10.1 asserts the chain
\begin{align}
  -\bigl|\log\bigl(q_\ell^{\varphi(y_0)}\bigr)\bigr| &= -\bigl|\log\bigl(q_\ell^{y_0}\bigr)\bigr|
  \tag{E1}\\[2pt]
  -\bigl|\log\bigl(q_\ell^{y_0}\bigr)\bigr|
  &\le -\tfrac{1}{\ell^*}\bigl|\log \Vol\bigl(\mathrm{hull}(\Thet^{I_M,\varphi(y_0)})\bigr)\bigr|
  \;=\; -\tfrac{1}{\ell^*}\bigl|\log \Vol\bigl(\mathrm{hull}(\Thet^{I,y_0})\bigr)\bigr|
  \;\le\; C_\Theta \bigl|\log\bigl(q_\ell^{y_0}\bigr)\bigr|
  \tag{E2}
\end{align}
\anchor{ATS~IV}{3180--3205}, where the first inequality is attributed to Mochizuki~\cite{Mochizuki2021c}
via the corresponding corollary of ATS~III. Our concern is
the middle equality of (E2) and the equality (E1) --- not the inequalities, which
we do not examine.

\subsection{E1: \texorpdfstring{$q_\ell$}{q\_l} is a classical, \texorpdfstring{$y$}{y}-free quantity}

The $q$-divisor of ATS~IV \S\,4.4 is defined by

\begin{quote}\small
``let $q_w$ at $C_{M_w}$ be a Tate parameter (\emph{well defined up to
multiplication by a local unit and its order at $w$}). \ldots\
$q_M = \sum_{w \in V^{\mathrm{odd},ss}_M} \ordd_w(q_w) \cdot w$''
\hfill\anchor{ATS~IV}{1831--1846}
\end{quote}

The definition chain used by the theorem can be displayed compactly as
\[
 q_M=\sum_{w\in V_M^{\mathrm{odd},ss}}\ordd_w(q_w)\,w,
 \qquad
 \log(q_M)=\frac{1}{[M:\mathbb{Q}]}\deg(q_M)
 =\frac{1}{[M:\mathbb{Q}]}
   \sum_w \ordd_w(q_w)\log w,
\]
followed, immediately before Theorem~6.10.1, by
$q_\ell^{y_0}=\tfrac{1}{2\ell}\log(q)$ computed in the $y_0$ frame
\anchor{ATS~IV}{1831--1876, 2523--2528, 3180--3189}. Here $q$ is the
\S\,5.4 shorthand for $q_{C/M}$, hence for the relevant \S\,4.4 divisor
$q_M$; correspondingly, $\log(q)$ denotes the normalized arithmetic degree
$\log(q_M)$, not a new object. Thus the superscript records the
computational frame; the displayed scalar factors through the classical divisor
and degree data.

so that $\ordd_w(q_w)$ is the classical discrete-valuation order in the local
field $M_w$; the parenthesis declares it a well-defined invariant of the
definition. The proof of Proposition~4.4.4 is correspondingly classical
throughout: uniformisers $\pi_w, \pi_v$, $q_v = \pi_v^\alpha$,
$\ordd_w(q_w) = e_{w\mid v}\ordd_v(q_v)$, $\log w = f_{w\mid v}\log v$
\anchor{ATS~IV}{1852--1876}. \emph{No $y$, no untilt and no holomorphoid occurs in
it.} The $q$-divisor is a $y$-free classical object of the pair $(C_M, M)$.

This settles E1 \emph{independently of the trilemma}, and the point is worth
making explicitly, because it is the one verdict here that needs no branch. The
quantities entering $q_\ell$ are integer orders $\ordd_w(q_w)$ and residue-field
cardinalities. Neither reads an absolute value supplied by an arithmeticoid, and
re-scaling a valuation by a positive exponent leaves the associated order
unchanged --- a re-normalization in any of the three senses of
Section~\ref{sec:trilemma} therefore cannot move $q_\ell$. So whichever branch is
intended, $q_\ell^{y_0}$ and $q_\ell^{\varphi(y_0)}$ compute the \emph{same
classical number}; the superscript records the ambient computational frame, but
does not change the value. Under branch~(a) the same conclusion can be reached via
N1$'$ --- both normalized copies read number-field elements in the same standard
scale --- but that route is not needed, and using it as the reason would make an
unconditional verdict look conditional. Direct visual readback of the ATS~IV
source PDF, p.~66, also shows absolute-value bars on \emph{both} sides of E1. The
one-sided version in the text extract was an extraction artefact. Hence no
additional assumption such as $\log(q_\ell)\ge 0$ is required: E1 is
definitionally true under the stated identification of $q_\ell$ with the
\S\,4.4 $q$-divisor.

\subsection{The apparent contradiction, and its resolution}

At first sight this collides with an explicit statement of the same paper: ``the
absolute values of Tate parameters rise under global Frobenius $\phi$''
\anchor{ATS~IV}{567--569}. If Tate parameters rise, how can a quantity built from
them be invariant?

The resolution is that two \emph{different objects} are involved. First, the
classical $q$-divisor of \S\,4.4, which is a combination of $\ordd_w$-values and
is $y$-free --- this is what E1 speaks about. Second, the Tate idele and the
untilt Tate parameters, which are $z$-dependent; these are what
\S\,1.4(2), Proposition~4.5.12 and the fiber statements of ATS~II speak about:

\begin{quote}\small
``Let $\arith(L)^{\mathrm{nor}}_z$ be the \emph{normalized} arithmeticoid \ldots\
the normalized arithmetic degree $\log(TI_{C/\arith(L)_z})$ \emph{is a function of
$z$} \ldots\ one obtains a \emph{non-constant} function `normalized degree of the
Tate divisor' $\Ynum \to \mathbb{R}$'' \hfill\anchor{ATS~IV}{2055--2075}
\end{quote}

Note that this non-constancy is asserted \emph{for the normalized} arithmeticoid:
normalization does not remove it, because it does not live on the number-field
level at all.

The same distinction disposes of a trap that the formulation phase had marked as
open. Joshi does document a holomorphoid dependence --- ``an important point of
the theory is that $\mathrm{Tate}(C/M)$ depends on the choice of a holomorphoid
$\hol(C/L)_y$ and one should write $\mathrm{Tate}(C/M)_y$ to indicate this
choice'' \anchor{ATS~IV}{681--684} --- but he documents it for
$\mathrm{Tate}(C/M)$, and he introduces $q$ precisely because $\mathrm{Tate}$
lacks a stability property:

\begin{quote}\small
``The divisors $\mathrm{Tate}(C/M)$, $\mathrm{Tate}_2(C/M)$ are not stable under
passage to field extensions \ldots\ so one replaces $\mathrm{Tate}(C/M)_2$ by the
normalized arithmetic degree of another divisor $q = q(C/M)$ (\S\,4.4), which is
similar to the divisor $\mathrm{Tate}$, but which receives contributions only from
primes of $M$ lying over a prime \ldots\ of semi-stable reduction \ldots\ By
Proposition~4.4.4, this new divisor $q$ is stable under base-field extensions.''
\hfill\anchor{ATS~IV}{698--705}
\end{quote}

Transferring the documented $y$-dependence of $\mathrm{Tate}$ to $q$ would be a
fallacy of object substitution; the index sets alone differ. What resolves the
matter is not a compensation mechanism but the observation that the $y$-variance
lives in a different formalization.

\subsection{E2: family stability and the open indexed equality}
The source PDF (ATS~IV v2, p.~66) equates \emph{absolute values} of
log-volumes. The former transcription without those bars asserted a stronger
statement. From $|a|=|b|$ one cannot infer $a=b$ without a common sign
binding; no such binding has been established here for the exact indexed hulls
and volume conventions in E2.

\paragraph{The source assertion: stability of the collated set.}
ATS~IV describes an adelic set which collates tuples of Tate parameters from
holomorphoids and says it is stable under iterates of the global Frobenius
\anchor{ATS~IV}{506--515}. This is evidence level~S, a source assertion.
It identifies a candidate family-level carrier; it does not identify the
$I_M,\varphi(y_0)$ and $I,y_0$ subsets, their ambient algebras, or their
measures. We have not traced those bindings.

\paragraph{Named transport: a sufficient stronger route, still open.}
Write $T:B_{I_M,\varphi(y_0)}\longrightarrow B_{I,y_0}$ for a proposed map.
One possible sufficient route would establish the ambient algebra and integral
structures, an isometry, the set identity
$T(\widetilde\Theta^{I_M,\varphi(y_0)})=\widetilde\Theta^{I,y_0}$,
hull compatibility
$T(\operatorname{hull}(S))=\operatorname{hull}(T(S))$,
and preservation of the relevant measures and degree weights.
These conditions would yield equality of the corresponding log-volumes,
hence their absolute values. They are stronger than the source's
absolute-value equality; they have not been shown necessary.
The map and each compatibility remain open. No replacement measure proof is
claimed from the generic hull counterexample or from a special algebra
decomposition.

\paragraph{The scale identification is only one layer.}
Under branch~(a), number-field elements have the standard absolute values in
each copy. Under branch~(b), the product formula restricts the conversion
multipliers to a global scalar, but does not bind that scalar across frames.
Neither observation proves E2 or transports normalized Haar measure.
Under branch~(c), even the element-scale identification requires its own
selection rule. The measure obligations of Section~\ref{sec:measure} are
open in all three branches.

\subsection{Where the work sits: the asymmetry of the two bounds}

A third textual observation completes the picture. The two bounds of the chain are
not computed in the same frame:

\begin{quote}\small
``by the virtue of stability under the global Frobenius morphism, one has the
freedom of the choice of an arithmeticoid in the computing the \emph{lower bound}
\ldots\ [it] is \emph{independent of the choice} \ldots\ the lower bound will be
computed using $y_0' = \varphi(y_0)$ while the upper bound will be computed using
$y_0$ (\emph{the upper bound is not independent of this choice}).''
\hfill\anchor{ATS~IV}{733--742}
\end{quote}

The source also notes that a prior release omitted this Frobenius shift
\anchor{ATS~IV}{755}. This locates a candidate role for the choice of frame;
it does not settle the equality. Branch~(a) makes the element-scale
identification definitional, while the shift, indexed sets, hulls and measures
still need compatible bindings. Branch~(c) additionally leaves the
selection-rule question. The present assessment does not verify either
inequality or certify that family stability supplies the missing indexed step.

\subsection{Classification table}

Table~\ref{tab:terms} records the classification. Two rows are deliberately not
forced into one of the invariance classes: $\Thet$ as a set is not a
scalar-valued quantity to be classified as invariant or not, and
$\log \Vol(\mathrm{hull}(\Thet))$ is the composite whose analysis is the subject
of Section~\ref{sec:measure}.

\begin{table}[htbp]
\centering
\caption{Term classification of Theorem~6.10.1 with assigned Evidence Levels. Class~T: definitional/trivially
invariant (factors over $\varphi$-fixed data). Class~N: demonstrably not
invariant (reads untilt data). Evidence levels: \textbf{D} (definitional), \textbf{S} (source-asserted), \textbf{T} (transport-verified), \textbf{O} (open or conditional). The classification is independent of the trilemma of
Section~\ref{sec:trilemma}; in the last row, indexed transport and cross-frame
measure compatibility remain open in every branch (Table~\ref{tab:branches}). Class~T and evidence level~T are distinct labels.}
\label{tab:terms}
\small
\begin{tabular}{@{}P{0.24\textwidth}P{0.20\textwidth}P{0.06\textwidth}P{0.11\textwidth}P{0.24\textwidth}@{}}
\toprule
Term & Factors over & Class & \shortstack{Evidence\\level} & Behavior \\
\midrule
$\ordd_w(q_w)$ (\S\,4.4)
  & discrete-valuation structure of $M_w$ ($\varphi$-fixed)
  & T & \textbf{D} & invariant \\[3pt]
$\log w$, $[M:\mathbb{Q}]$, $e_{w\mid v}$, $f_{w\mid v}$
  & residue and degree data ($\varphi$-fixed)
  & T & \textbf{D} & invariant \\[3pt]
$q_\ell = \tfrac{1}{2\ell}\log(q)$
  & only the above
  & \textbf{T} & \textbf{D} & \textbf{E1 definitional}; source PDF has absolute values on both sides \\[3pt]
Tate idele degree $\log(TI_{C/\arith(L)_z})$
  & untilt valuations
  & \textbf{N} & \textbf{S} & non-constant (Prop.\ 4.5.12) \\[3pt]
Kummer embedding $\mathbb{Z}_p(1) \to B$
  & choice of untilt $y_n$
  & N & \textbf{S} & differs by a $p$-multiple (Thm.\ 10.20.1(2)) \\[3pt]
$\Thet_{\mathrm{Mochizuki}}$ as a set
  & the whole family (collated)
  & --- & \textbf{S} & $\varphi$-stable by construction \\[3pt]
$\log \Vol(\mathrm{hull}(\Thet))$
  & set $+$ measure convention
  & E2 & \textbf{O} & Absolute-value E2: indexed transport and measure compatibility open in every branch \\
\bottomrule
\end{tabular}
\end{table}

Criterion~K has illustrative examples in the classical $q$-divisor and
the untilt-sensitive terms. These examples do not make it a general
invariance theorem and do not classify the algebraic hull or the weighted
cross-frame log-volume as transport-verified.

% =============================================================================
\section{Measure conventions and open cross-frame obligations}
\label{sec:measure}

\paragraph{Absolute value, Haar measure and multiplication module.}
For each place and frame, distinguish the absolute value $|\cdot|_{y,v}$,
the additive Haar measure $m_{y,v}$ normalized by
$m_{y,v}(\mathcal O_{L'_w})=1$, and its multiplication module
\[
 \Delta_{y,v}(\alpha)=\frac{m_{y,v}(\alpha A)}{m_{y,v}(A)}
\]
for a measurable set $A$ of positive finite measure and $\alpha\ne0$.
The equality $\Delta_{y,v}(\alpha)=|\alpha|_{y,v}$ requires a compatible
absolute-value convention. On $\mathbb Q_p$, normalization forces
\[
 m(\mathbb Z_p)=1,\quad m(p\mathbb Z_p)=p^{-1},\quad |p|_p^c=p^{-c}.
\]
The second identity follows by partitioning $\mathbb Z_p$ into its $p$
cosets modulo $p\mathbb Z_p$. Changing a valuation exponent does not change
this normalized additive Haar measure. The scalar in
Lemma~\ref{lem:rigidity} therefore provides no measure-transfer proof.

\paragraph{Local volume conventions and degree weights.}
ATS~III defines normalized Haar measure on the number-field completions
by $\Vol_{L'_w}(\mathcal O_{L'_w})=1$
\anchor{ATS~III}{6562--6572} and relates suitable fractional-ideal coset
volumes to local absolute values
\anchor{ATS~III}{6599--6620,6663--6676}.
It specifies reciprocal local degrees as weights
\anchor{ATS~III}{6678--6688}.
These are statements about specified fields, integral structures and
measure conventions. Their numerical ingredients are classical; they do not
identify the ambient tensor algebras, exact theta sets, hulls, weighted
components or measures across E2 frames. Standard element scales under
branch~(a), or a common valuation exponent under branch~(b), are insufficient.
Branch~(c) additionally leaves the element-scale choice open.

\paragraph{Algebraic hull need not equal convex closure.}
ATS~III v4, \S\S\,9.10.7--9.10.8 (pp.~125--127), defines the holomorphic
hull through integer rings of field components of the tensor algebra.
The former generic identification with convex closure was incorrect.
Take $E/\mathbb Q_p$ unramified quadratic, $E_1=E$,
$E_2=\mathbb Q_p$, the identity algebra map, and $U=\mathbb Z_p\cdot1$.
The set is compact and already $\mathbb Z_p$-convex. The smallest hull
of the prescribed form $\alpha\mathcal O_E$ containing $1$ is
$\mathcal O_E$: if $1\in\alpha\mathcal O_E$, then
$\alpha\mathcal O_E\supseteq\mathcal O_E$, with equality attained for
$\alpha=1$. Thus
\[
 \operatorname{conv}(U)=\mathbb Z_p\cdot1
 \subsetneq\mathcal O_E=\operatorname{hull}(U).
\]
Strictness follows from the $\mathbb Z_p$-ranks one and two.
Convexity of a hull containing $U$ gives only
$\operatorname{conv}(U)\subseteq\operatorname{hull}(U)$ in general.
Admissibility of this $U$ as a specialized Mochizuki tensor product region
has not been established. The example refutes our generic identification,
not Proposition~9.10.8.1(3), Theorem~9.11.1, or $abc$.
A proposed replacement through decomposition of finite étale algebras
still needs the actual algebra maps, integral structures, sets, weights
and indexed transport; it is not a verified positive measure argument.

\paragraph{Revised conclusion.}
The former assertion that the measure side factors without remainder over
number-field data, or is pinned under branch~(a) and up to a scalar under
branch~(b), is withdrawn. Cross-frame measure compatibility remains open
in every branch. Identifying element scales does not close indexed E2.

\begin{table}[htbp]
\centering
\caption{Components after correction. D: definitional; S: source-asserted;
O: open or conditional.}
\label{tab:resolution}
\small
\begin{tabular}{@{}P{0.23\textwidth}P{0.27\textwidth}P{0.08\textwidth}P{0.29\textwidth}@{}}
\toprule
Component & Basis & Level & Status \\
\midrule
Element scales & Place-wise standard reading & O
 & Definitional under (a); global scalar under (b); open selection under (c) \\[3pt]
E1 & Classical $q$-divisor & D
 & Under the stated identification; absolute values on both sides \\[3pt]
Collated family & Frobenius-stability assertion & S
 & Candidate carrier; indexed E2 not established \\[3pt]
Indexed E2 & Absolute values of log-volumes & O
 & Open in every branch; unsigned equality not inferred \\[3pt]
Measure and hull & Haar modules, integral structures, algebraic hull, degree weights & O
 & Compatibility open in every branch; generic hull equality withdrawn \\[3pt]
Selection rule & Coordinate $\lambda_y$ & O
 & Open correspondence question \\[3pt]
Local comparison & ATS~IV Remark~6.10.2 & S
 & Expressly precluded \\
\bottomrule
\end{tabular}
\end{table}

% =============================================================================
\section{Where the trilemma comes from: the height passage}
\label{sec:fneu}
% =============================================================================

The trilemma of Section~\ref{sec:trilemma} was not part of the question this paper
set out to answer. It emerged from a passage encountered while reading, and this
section records how, because the guiding principle of the program is to document
what was seen rather than to defend a position.

\begin{findingF}
The height passages of ATS~II(1/2) \S\S\,8.2--8.3 are indicators of a
possible tension under the place-wise reading. The claimed non-constancy,
its proof using a local valuation, and its domain $\mathbb P^n(\overline L)$
must be checked separately. A selection rule for $\lambda_y$ may clarify
the reading, but this paper has not settled the domain or proof questions.
F-NEU-1 remains open.
\end{findingF}

The definition uses the \emph{normalized} valuations explicitly: ``Choose the
standard normalization \ldots\ work with $\arith(L)^{\mathrm{nor}}_y$ \ldots\ work
with the normalization coordinate $\lambda_y$'', giving
$h(P) = \sum_v \max_j \log\bigl(|x_j|^{\lambda_v}_{K_{y_v}}\bigr)$
\anchor{ATS~II(1/2)}{2947--2960}; and Lemma~8.2.3 derives the independence of the
height from the choice of homogeneous coordinates precisely from the fact ``that
we work with a normalized arithmeticoid'' \anchor{ATS~II(1/2)}{2982--2986}. Under
a place-wise reading of (5.3.3) these normalized values agree, on elements of
$L_v$, with the standard ones --- for \emph{every} $y$ --- so for
$P \in \mathbb{P}^n(L)$ the height would be the classical, $y$-independent height
in the sense of~\cite{BombieriGubler2006}. ATS~IV says as much about the
definition: ``except for explicating the dependence on the arithmeticoid, the
definition of height \ldots\ is the same as the one given in [Bombieri and Gubler,
2006, \S\,1.5.1]'' \anchor{ATS~IV}{548--550}.

Proposition~8.3.1, however, asserts that
``$h_{\arith(L)}$ depends \emph{non-trivially} on $\arith(L)$''
\anchor{ATS~II(1/2)}{2990--3004}, and its proof invokes the fact that
``$y \mapsto |p|_{K_y}$ is a non-constant function'' --- which is the
\emph{canonical}, unnormalized valuation function of Remark~2.7.1. Remark~8.3.2
adds that ``the normalization of $\arith(L)^{\mathrm{nor}}_{y_1}$ \emph{cannot be
carried over} to $\arith(L)^{\mathrm{nor}}_{y_2}$''
\anchor{ATS~II(1/2)}{3008--3015}, which closes the one easy escape --- keeping
$\lambda_{y_1}$ fixed while $y$ moves.

Three possible readings require separate checks. The definition in
ATS~II(1/2) \S\,8.2 and Lemma~8.2.3 uses a normalized height; its
relationship to canonical local valuations must not be settled by replacing
one object with another. In particular, Proposition~8.3.1 states
non-constancy of that normalized height on $\mathbb P^n(\overline L)$,
although its proof also uses an unnormalized local absolute value. The
statement, proof step and domain remain distinct review obligations.

The proposed restriction of $y$-dependence to $\overline L\setminus L$
cannot be rejected merely from ATS~IV:545: the original PDF v2, p.~13,
says $x\in\overline L^*$, not $x\in L^*$. Its application to integer
$abc$-triples is a separate passage \anchor{ATS~IV}{608--616}.
Independently, ATS~II(1/2) v12, Theorem~8.7.1 (pp.~54--55), really
uses $x\in L^*$ and asserts general strictness
\anchor{ATS~II(1/2)}{3210--3216}.
This is a distinct textual indication, not a substitute for the incorrect
ATS~IV quotation. The family assertion that normalization cannot be carried
over does not by itself produce strictness at a fixed point.
F-NEU-1 therefore remains an open correspondence question; none of these
observations establishes a global error in ATS.

What the passage does supply is the mechanism, and it is the mechanism that
generates the trilemma. In \S\,8.5 the movement of the height is derived from the
product-formula constraint:

\begin{quote}\small
``since one is working with normalized arithmeticoids $\arith(L)^{\mathrm{nor}}$,
the normalization constraint i.e.,\ the product formula (5.3.4) implies that
\emph{changes in $p$-adic valuations forced by the action of $n$ must be
compensated by changes at other primes} so that the normalization constraint
\ldots\ continues to hold.'' \hfill\anchor{ATS~II(1/2)}{3162--3193}
\end{quote}

Compensation at other places presupposes freedom at those places. So the operative
meaning of $\arith(L)^{\mathrm{nor}}$ in this passage is not ``bound place by
place to the Artin--Whaples standard'' but ``constrained, globally, by the product
formula'' --- and then the question is which rule picks the coordinate out of the
admissible set. That is branch (b) or (c) of Section~\ref{sec:trilemma}, and
Lemma~\ref{lem:rigidity} says that (b) does not deliver what the passage needs,
since the freedom the product formula leaves is a single global scalar and cannot
be place-selective.

\paragraph{The resulting question, and its status.}
The precise question we would put to the author is therefore this: \emph{which
rule selects the normalization coordinate $\lambda_y$ for a moved arithmeticoid?}
Concretely --- is the normalized arithmeticoid the one whose induced absolute
values on $L_v$ are the Artin--Whaples standard ones, as (5.3.3) suggests, or is
$(\cdot)^{\mathrm{nor}}$ a constraint on a choice, for which the product formula
alone is not sufficient to pin down $\lambda_y$?

We make no error claim, against Proposition~8.3.1 or against anything else. The
resolving passage may well stand in material we have not read, or follow from a
reading of the definition that the author would supply at once. The two textual
observations above --- the place-wise appearance of (5.3.3), and the
place-selective use of the normalization in \S\,8.5, in ATS~III Corollary~4.2.2.4
and in \S\,1.5 of ATS~IV --- are reported as observations. What has changed
relative to the first version of this assessment is only this: the tension cannot
be set aside as marginal, because the verdicts of
Sections~\ref{sec:trichotomy}--\ref{sec:measure} depend on how it is resolved.
That dependence is what Table~\ref{tab:branches} records. The reading of the
sources underlying this section has been carried out in more detail in the
companion verification of this program (Parts~V and~VI, in preparation).

% =============================================================================
\section{Assessment and consequences}
\label{sec:assessment}
% =============================================================================

\subsection{For the ledger method}

The original ledger located a question at the phrase declaring the
normalization identification a tautology. Under the place-wise reading,
that phrase indeed identifies number-field element scales definitionally.
This limited conclusion does not discharge E2: the source's stability
assertion concerns a collated family, whereas the theorem evaluates two
indexed algebraic hulls with weighted volume conventions.

The corrected audit separates these duties. It reports the family assertion
at evidence level~S and keeps the indexed equality, sign binding and
cross-frame measure compatibility at level~O. The generic hull identification
used in the earlier positive measure argument is withdrawn.
Under branch~(c), selection of the normalization coordinate is a further
open duty. A formally acceptable ledger or a source-citation check is not
a proof of those duties and does not strengthen an unverified bound.

The method's useful output here is the correction of the actual source
readings and a more precise set of open obligations. It is not a completed
bridge certificate or acceptance of the entire ATS argument.

\begin{table}[htbp]
\centering
\caption{Selected Gap-Series papers (non-exhaustive crosswalk): Target claims, primary diagnostic instruments, and core normalization/transport verdicts.}
\label{tab:gap_series_crosswalk}
\small
\setlength{\tabcolsep}{3.5pt}
\begin{tabular}{@{}P{0.14\textwidth}P{0.19\textwidth}P{0.20\textwidth}P{0.27\textwidth}P{0.12\textwidth}@{}}
\toprule
Part / Paper & Target Claim & Primary Instrument & Core Verdict & Status \\
\midrule
\textbf{Part~I}\newline (Method)
  & Inter-universal transport
  & Bridge-Contract Criterion (C1--C7)
  & Foundational methodology; specifies explicit verification duties
  & Published \\[3pt]
\textbf{Part~IV}\newline (This Paper)
  & Global normalization in ATS~IV Thm~6.10.1
  & Invariance classification \& Trilemma taxonomy
  & Identification is definitional under (a); transport of $\Thet$-hull untraced
  & Corrective draft v1.2 \\[3pt]
\textbf{Part~IV-B}\newline (BrXi App.)
  & Self-application of $\mathrm{Br}_{\Xi}$
  & Sub-bridge ledger audit
  & Scope restriction; prevents cyclical domain inheritance across boundaries
  & Draft v0.1 \\[3pt]
\textbf{Part~V}\newline (Substance)
  & Normalization weights in ATS
  & All-place rigidity classification
  & Excludes non-trivial place-selective scaling under global Artin--Whaples
  & Draft v0.1 \\[3pt]
\textbf{Part~VI}\newline (Repair)
  & Repair program for ATS chain
  & Effective-weight test \& Case taxonomy
  & Proves all-place rigidity for product-formula weights; tests height chain
  & Draft v0.2 \\
\bottomrule
\end{tabular}
\end{table}

\subsection{For global-only mechanisms: a cautious outlook}
\label{sec:b2outlook}

The structural pattern of Theorem~6.10.1 --- ``no comparison place by place, pinned
globally'' --- is not peculiar to this construction. Three carriers of that pattern
appear in the material, two of them from Joshi's own text.

The prototype is the product formula itself~\cite{ArtinWhaples1945}. The identity
$\prod_v |x|_{K_v} = 1$ says nothing whatever about the individual $|x|_v$, and is
nevertheless the load-bearing global constraint of the entire construction. The
second carrier is Joshi's own second example within the same proof passage:
``$A' = C \cdot A$. \emph{This is a global equation, which only appears after one
sums over all primes}'' \anchor{ATS~IV}{834--836}. That two such statements sit
next to each other in one argument suggests a recurring pattern rather than an
idiosyncrasy of the middle equality. The third is the pattern examined in this
paper, now available as a worked specimen: the hyperplane construction plus
Remark~6.10.2 give an instance in which the global pinning and the local
incomparability are both explicit and both sourced.

A working hypothesis of the present program --- recorded in its own candidate
register, and stated here with the caution it requires --- is that such statements
live on the collated family rather than at the individual place, and that this
gives them a structural silhouette in common with statements available only after
summation over all places, Weil positivity among them. We state this as a
hypothesis about the \emph{shape of arguments}. It is emphatically \emph{not} a
mathematical connection, and \emph{not} a claim about the Riemann hypothesis or
about any consequence of one. The corresponding work package carries an explicit
abort clause: should the collected carriers turn out to share nothing beyond the
observation itself, the analogy is to be recorded as a negative result and the
register entry downgraded --- not quietly carried forward. A reverse-direction
question in the vicinity of~\cite{GranvilleStark2000} is noted in the project's
concept document and is expressly not taken up here.

\subsection{Source corrections and open obligations}
\label{sec:residues}

\paragraph{Three corrected readings.}
The source-PDF domain in ATS~IV v2, p.~13, is $\overline L^*$.
ATS~II(1/2) Theorem~8.7.1 separately uses $L^*$; the two anchors
are not interchangeable (Section~\ref{sec:fneu}).
ATS~IV v2, p.~66, places absolute-value bars around both log-volume
terms in E2. No common sign of the exact indexed evaluations is proved
here, so the former unsigned equality is withdrawn.
ATS~III v4, pp.~125--127, uses an algebraic hull; its generic equality
with convex closure is refuted by the explicit example in
Section~\ref{sec:measure}. This does not decide the special tensor-region
claims. The earlier E1 sign discrepancy remains separately resolved as an
extraction artefact: the source has absolute-value bars on both E1 sides.

\paragraph{Open indexed equality and measure transfer.}
Family stability is source-asserted. A map between the exact indexed
ambient algebras, its integral structures, theta-set images, hull action,
measures and weights has not been traced. The named transport is only
a sufficient stronger route to E2; it is not a necessary condition derived
from the source. Neither a norm power nor a generic hull identity supplies
the absent measure argument. These obligations remain open in every branch.

\paragraph{Selection and order questions.}
The rule selecting $\lambda_y$ remains open. Lemma~\ref{lem:rigidity}
constrains only multipliers of standard valuations satisfying the product
formula; it supplies no place-selective rule or Haar-measure transport.
D-1 remains conditional on (F1)--(F2). F-NEU-1 and the strong order
question L-strong remain open. Separate follow-up audits are needed before
positive E2 or measure statements from companion Parts~V and~VI can be
imported as verified results.

\subsection{Limitations}
\label{sec:limitations}

This subsection states what the paper does not do, in the terms in which the
underlying working notes stated it.

\begin{enumerate}
\item \textbf{No proof was followed.} Neither Theorem~6.10.1 of ATS~IV, nor the
  corollary of ATS~III on which its first inequality relies, nor Mochizuki's
  Corollary~3.12~\cite{Mochizuki2021c} was verified line by line. No estimate was
  recomputed and no bound was checked. What was classified is the \emph{terms and
  the structure of justification}, not the correctness of the inequalities.
\item \textbf{D-1 is our own derivation and is conditional}, not a theorem of the
  source: elementary algebra plus a proportionality argument, assembled from
  statements quoted verbatim. It holds under the assumptions (F1)--(F2) stated with
  it --- the canonical-fiber assumption for the standard point, and the transfer of
  the coordinate transformation to two finite places, which is read off the
  quotation rather than computed there. The archimedean case is likewise unchecked
  and, for D-1, not needed.
\item \textbf{Element-scale verdicts are branch-conditional.}
  Definitional identification and the vacuity of N3 apply under branch~(a);
  branch~(b) leaves a global conversion scalar and branch~(c) its selection rule.
  Cross-frame measure pinning is not proved under any branch and its former
  positive classification is withdrawn. D-1 requires (F1)--(F2); E1 retains
  its stated identification; family stability is source-asserted, not an
  indexed hull or volume proof.
\item \textbf{The trilemma is a definitional gap, not an error claim.} We did not
  find a selection rule for $\lambda_y$ in the passages we read; we do not claim
  that none exists, and we make no claim about the correctness of ATS~III
  Corollary~4.2.2.4 or of anything resting on it.
\item \textbf{Lemma~\ref{lem:rigidity} is stated with a deliberate reference
  restriction.} The $\mathbb{Q}$-case is proved here in full; the general case
  invokes Dirichlet's unit theorem and the finiteness of the class number as
  classical results without a section number, because we verified none. In
  particular we do not attribute the lemma to~\cite{ArtinWhaples1945}, whose full
  text we could not obtain.
\item \textbf{F-NEU-1 remains open.} The corrected ATS~IV domain
  $\overline L^*$ and the separate ATS~II(1/2) statement on $L^*$ do not
  adjudicate the entire height argument or establish an error in ATS.
\item \textbf{The classification of E1 as definitional} holds modulo the reading
  that $q_\ell$ in Theorem~6.10.1 is the
  \S\,4.4 $q$-divisor. The definition given immediately before the theorem
  supports that reading, but the identification is ours. The source PDF itself
  displays absolute-value bars on both sides of E1.
\item \textbf{The measure analysis classifies conventions, not estimates.} The
  correctness of the volume estimate itself --- the product bound, the
  $\ell^*$-factors, the translation to Mochizuki's formulation --- was not
  examined.
\item \textbf{The strong form of the question is undeveloped.} The order
  compatibility question (L-strong) of Section~\ref{sec:equivariance} names a
  structure but no candidate class. This is the largest substantive arrear.
\item \textbf{The invariance catalogue is a starting inventory}, not a
  completeness claim. Arakelov and adelic degrees in particular remain unfilled,
  because the material examined says nothing usable about them and no reference
  was guessed.
\item \textbf{Textual basis and visual readback.} Quotations use
  \texttt{pdftotext} extracts; line references are finding aids, not pagination.
  The source-PDF pages for equation (5.3.3), Theorem~5.10.1(3)--(4),
  Theorem~10.20.1(3), Definition~4.4.2/formula (4.4.5), and
  Theorem~6.10.1 were visually checked. In particular, exponent direction and the
  two pairs of E1 absolute-value bars are load-bearing. No blanket OCR-robustness
  claim is made for other passages.
\item \textbf{No claim upgrade.} Nothing in this paper is a statement about the
  correctness of Joshi's $abc$ argument, about Mochizuki's theory~\cite{Mochizuki2021a,Mochizuki2021b,Mochizuki2021c,Mochizuki2021d}, about the
  objections of~\cite{ScholzeStix2018}, or about the Riemann hypothesis. It is the
  clarification of one definitional question and of its consequence for one
  conjectured lemma.
\end{enumerate}

% =============================================================================
% AI DISCLOSURE -- pipeline standard, included from the canonical template
% (_templates/disclosure/ai_disclosure_STANDARD.tex, bilingual, carries its own
% \section*). The template names computational scripts and runs; this paper has
% none, so the restriction paragraph below adapts it downwards, which is what the
% template's own instruction ("Abweichen nach unten") requires. Do not silently
% drop the restriction: without it the disclosure would claim work not done.
% =============================================================================
\phantomsection
\addcontentsline{toc}{section}{AI Disclosure}
% =============================================================================
% AI DISCLOSURE -- PIPELINE STANDARD, ENGLISH VERSION (since 2026-08-17, T-20260817-02)
% =============================================================================
% Language-pure version for English-only manuscripts. Text copied verbatim
% from the English block of ai_disclosure_STANDARD.tex (no content change) --
% only the heading is now language-pure and the German block is dropped. The
% bilingual file ai_disclosure_STANDARD.tex remains unchanged (existing
% references may keep pointing to it).
%
% DEVIATING DOWNWARD: If a paper genuinely used substantially less AI, do NOT
% switch to a different file -- weaken the wording here instead, e.g. replace
% "extensively in all phases" with the phases that actually applied, and
% strike whatever did not happen from the quality-assurance list. Principle:
% state honestly what happened -- do not hedge.
%
% Usage (English-only manuscripts only):
%   \input{../../_templates/disclosure/ai_disclosure_STANDARD_en}
% =============================================================================

\section*{AI Disclosure}

\noindent
AI systems were used extensively in all phases of this work: conception
and research discussion, literature and source work, mathematical analysis
and proof development, creation and checking of the computational scripts,
execution and evaluation of the computational runs, text, translation, and
typesetting.

\medskip

\noindent
Quality assurance: use of several independent models from different
providers with counter-checks and adversarial AI reviews; traceable,
versioned computational scripts with stored result files and checksums;
continuously maintained proof and register documents; strict claim hygiene
(assertions only with documented evidence, addenda instead of silent
corrections); changelogs across all versions.


\medskip

\noindent
\emph{Restriction for this paper:} this work involved no computational scripts and
no computational runs; the corresponding items of the standard text (creation and
checking of the computational scripts, execution and evaluation of the
computational runs, versioned scripts with stored result files and checksums) do
not apply here.

% =============================================================================
% BIBLIOGRAPHY
% Every entry is provisional and must be checked against the original before any
% upload or submission. No entry was included that is not used in the text or
% expressly referenced by the working notes.
% =============================================================================
\begin{thebibliography}{99}
\addcontentsline{toc}{section}{References}
\small

% verified 2026-08-13 — Quellencheck (WebSearch bzw. lokale Extrakte), Beleg: _sources/QUELLENCHECK_BIBITEMS_2026-08-13.md
\bibitem{JoshiATSIIhalf}
K.~Joshi.
\newblock \emph{Construction of Arithmetic Teichmüller Spaces II(1/2):
Deformations of Number Fields}.
\newblock arXiv:2305.10398, version v12 (24 February 2025). Cited here as
ATS~II(1/2); ``[Joshi, 2023a]'' in the author's own bibliographies.

% verified 2026-08-13 — Quellencheck (WebSearch bzw. lokale Extrakte), Beleg: _sources/QUELLENCHECK_BIBITEMS_2026-08-13.md
\bibitem{JoshiATSII}
K.~Joshi.
\newblock \emph{Construction of Arithmetic Teichmüller Spaces II: Proof of a
local prototype of Mochizuki's Corollary~3.12}.
\newblock arXiv:2303.01662, version v3 (24 February 2025). Cited here as ATS~II;
``[Joshi, 2023b]'' in the author's own bibliographies.

% verified 2026-08-13 against local extract header (arXiv:2401.13508v4)
\bibitem{JoshiATSIII}
K.~Joshi.
\newblock \emph{Construction of Arithmetic Teichmüller Spaces III: A `Rosetta
Stone' and a proof of Mochizuki's Corollary~3.12}.
\newblock arXiv:2401.13508, version v4 (24 February 2025). Cited here as
ATS~III; ``[Joshi, 2024c]'' in the author's own bibliographies.

% verified 2026-08-13 against local extract header (arXiv:2403.10430v2)
\bibitem{JoshiATSIV}
K.~Joshi.
\newblock \emph{Construction of Arithmetic Teichmüller Spaces IV: Proof of the
abc-conjecture}. Preliminary version for comments.
\newblock arXiv:2403.10430, version v2 (24 February 2025). Cited here as ATS~IV.

% verified 2026-08-13 — Quellencheck (WebSearch bzw. lokale Extrakte), Beleg: _sources/QUELLENCHECK_BIBITEMS_2026-08-13.md
\bibitem{Mochizuki2021a}
S.~Mochizuki.
\newblock \emph{Inter-universal Teichmüller Theory I: Construction of Hodge
Theaters}.
\newblock Publ. Res. Inst. Math. Sci. \textbf{57} (2021), no.~1/2, 3--207.

% verified 2026-08-13 — Quellencheck (WebSearch bzw. lokale Extrakte), Beleg: _sources/QUELLENCHECK_BIBITEMS_2026-08-13.md
\bibitem{Mochizuki2021b}
S.~Mochizuki.
\newblock \emph{Inter-universal Teichmüller Theory II: Hodge--Arakelov-theoretic
Evaluation}.
\newblock Publ. Res. Inst. Math. Sci. \textbf{57} (2021), no.~1/2, 209--401.

% verified 2026-08-13 — Quellencheck (WebSearch bzw. lokale Extrakte), Beleg: _sources/QUELLENCHECK_BIBITEMS_2026-08-13.md
\bibitem{Mochizuki2021c}
S.~Mochizuki.
\newblock \emph{Inter-universal Teichmüller Theory III: Canonical Splittings of
the Log-theta-lattice}.
\newblock Publ. Res. Inst. Math. Sci. \textbf{57} (2021), no.~1/2, 403--626.
\newblock Corollary~3.12 and Remark~3.1.1 are the passages referenced through
Joshi's text.

% verified 2026-08-13 — Quellencheck (WebSearch bzw. lokale Extrakte), Beleg: _sources/QUELLENCHECK_BIBITEMS_2026-08-13.md
\bibitem{Mochizuki2021d}
S.~Mochizuki.
\newblock \emph{Inter-universal Teichmüller Theory IV: Log-volume Computations
and Set-theoretic Foundations}.
\newblock Publ. Res. Inst. Math. Sci. \textbf{57} (2021), no.~1/2, 627--723.

% verified 2026-08-13 — Quellencheck (WebSearch bzw. lokale Extrakte), Beleg: _sources/QUELLENCHECK_BIBITEMS_2026-08-13.md
\bibitem{ScholzeStix2018}
P.~Scholze and J.~Stix.
\newblock \emph{Why abc is still a conjecture}.
\newblock Manuscript, 2018. Available at
\texttt{https://www.math.uni-bonn.de/people/scholze/WhyABCisStillaConjecture.pdf}.

% verified 2026-08-13 — Quellencheck (WebSearch bzw. lokale Extrakte), Beleg: _sources/QUELLENCHECK_BIBITEMS_2026-08-13.md
\bibitem{ArtinWhaples1945}
E.~Artin and G.~Whaples.
\newblock \emph{Axiomatic characterization of fields by the product formula for
valuations}.
\newblock Bull. Amer. Math. Soc. \textbf{51} (1945), 469--492.

% v1.1: single new entry, admitted because Lemma~1 and Section 3.3(i) need a
% verifiable locator for the classical uniqueness statement. Fundstelle [V]:
% PDF eingesehen, Prop. 7.8 auf S. 108 im Wortlaut geprueft (2026-08-13).
% Milne's numbering is version-dependent -- version and page must stay quoted.
\bibitem{Milne2020}
J.~S.~Milne.
\newblock \emph{Algebraic Number Theory}.
\newblock Course notes, version 3.08 (19 July 2020). Available at
\texttt{https://www.jmilne.org/math/CourseNotes/ANT.pdf}. Cited by version, item
number and page, because the numbering differs between versions.

% verified 2026-08-13 — Quellencheck (WebSearch bzw. lokale Extrakte), Beleg: _sources/QUELLENCHECK_BIBITEMS_2026-08-13.md
\bibitem{FarguesFontaine2018}
L.~Fargues and J.-M.~Fontaine.
\newblock \emph{Courbes et fibr\'es vectoriels en th\'eorie de Hodge $p$-adique}.
\newblock Ast\'erisque \textbf{406}, Soci\'et\'e Math\'ematique de France, 2018.
Pr\'eface par Pierre Colmez, xiii+382~pp.

% verified 2026-08-13 — Quellencheck (WebSearch bzw. lokale Extrakte), Beleg: _sources/QUELLENCHECK_BIBITEMS_2026-08-13.md
\bibitem{BombieriGubler2006}
E.~Bombieri and W.~Gubler.
\newblock \emph{Heights in Diophantine Geometry}.
\newblock New Mathematical Monographs 4, Cambridge University Press, 2006.

% verified 2026-08-13 — Quellencheck (WebSearch bzw. lokale Extrakte), Beleg: _sources/QUELLENCHECK_BIBITEMS_2026-08-13.md
\bibitem{GranvilleStark2000}
A.~Granville and H.~M.~Stark.
\newblock \emph{ABC implies no ``Siegel zeros'' for $L$-functions of characters
with negative discriminant}.
\newblock Invent. Math. \textbf{139} (2000), 509--523.

\end{thebibliography}

\end{document}
